\section{Solutions}\label{sec:solutions_linear_equations}

% Solution by Gemini 3.7 Flash (High)
\begin{exercisesolution}[Solution to \cref{exc:converse_ero_solution_set}]
\label{sol:converse_ero_solution_set}
\exinfo{Solution by Gemini 3.7 Flash (High). Examines row equivalence vs solution set equality for inconsistent and non-square systems.}
The converse is \textbf{false} in general. Having identical solution sets does not imply that the augmented matrices are row equivalent.

\textbf{Counterexample 1 (Inconsistent Systems):}
Consider the two linear systems over $\mathbb{R}$ in two variables $(x_1, x_2)$:
\[
(S_1): \begin{cases} x_1 + 0x_2 = 0 \\ 0x_1 + 0x_2 = 1 \end{cases} \qquad\text{and}\qquad (S_2): \begin{cases} 0x_1 + x_2 = 0 \\ 0x_1 + 0x_2 = 1 \end{cases}
\]
Both systems are inconsistent, so their solution sets are both empty: $L(S_1) = \emptyset = L(S_2)$.
Their augmented matrices in row-reduced echelon form are:
\[
(A_1 \mid b_1) = \begin{pmatrix} 1 & 0 & \big| & 0 \\ 0 & 0 & \big| & 1 \end{pmatrix} \qquad\text{and}\qquad (A_2 \mid b_2) = \begin{pmatrix} 0 & 1 & \big| & 0 \\ 0 & 0 & \big| & 1 \end{pmatrix}.
\]
These two matrices are not row equivalent because the only non-zero row in the coefficient part of $(A_1 \mid b_1)$ is $(1, 0)$, whereas in $(A_2 \mid b_2)$ it is $(0, 1)$. Since both matrices are in row-reduced echelon form (\cref{def:row_reduced_echelon}), they represent different row spaces and cannot be transformed into one another by elementary row operations.

\textbf{Counterexample 2 (Different Number of Equations):}
The system $(x_1 = 0)$ with augmented matrix $(1 \mid 0) \in \M_{1 \times 2}(\mathbb{R})$ and the system $\begin{cases} x_1 = 0 \\ x_1 = 0 \end{cases}$ with augmented matrix $\begin{pmatrix} 1 & 0 \\ 1 & 0 \end{pmatrix} \in \M_{2 \times 2}(\mathbb{R})$ both have the unique solution $x_1 = 0$. However, matrices of different dimensions cannot be row equivalent.

\begin{remark}
If one restricts attention to \emph{consistent} linear systems with the same number of variables and equations whose coefficient matrices have the same dimensions, then the converse does hold: any consistent system has a unique reduced row echelon form completely determined by the affine subspace of solutions.
\end{remark}
\end{exercisesolution}

% Solution by Gemini 3.7 Flash (High)
\begin{exercisesolution}[Solution to \cref{exc:gauss_elimination_4vars}]
\label{sol:gauss_elimination_4vars}
\exinfo{Solution by Gemini 3.7 Flash (High). Fully works out Gaussian elimination into reduced row echelon form with 2 free variables.}
We write the augmented matrix of the system and apply elementary row operations:
\[
\begin{pmatrix}
1 & 2 & 2 & 1 & \big| & -1 \\
3 & 6 & 2 & 5 & \big| & 1
\end{pmatrix}
\xrightarrow{-3 \cdot E_1 + E_2 \to E_2}
\begin{pmatrix}
1 & 2 & 2 & 1 & \big| & -1 \\
0 & 0 & -4 & 2 & \big| & 4
\end{pmatrix}
\]
Divide the second row by $-4$:
\[
\xrightarrow{-\frac{1}{4} \cdot E_2 \to E_2}
\begin{pmatrix}
1 & 2 & 2 & 1 & \big| & -1 \\
0 & 0 & 1 & -\frac{1}{2} & \big| & -1
\end{pmatrix}
\]
Now clear the entry above the second pivot by $-2 \cdot E_2 + E_1 \to E_1$:
\[
\xrightarrow{-2 \cdot E_2 + E_1 \to E_1}
\begin{pmatrix}
1 & 2 & 0 & 2 & \big| & 1 \\
0 & 0 & 1 & -\frac{1}{2} & \big| & -1
\end{pmatrix}
\]
The matrix is now in row-reduced echelon form. 
\begin{itemize}
    \item \textbf{Pivot variables:} $x$ (column 1) and $z$ (column 3).
    \item \textbf{Free variables:} $y$ (column 2) and $w$ (column 4).
\end{itemize}
Let $y = s \in \mathbb{R}$ and $w = t \in \mathbb{R}$ be arbitrary parameters. Expressing the pivot variables in terms of the free parameters gives:
\[
\begin{cases}
x = 1 - 2s - 2t, \\
z = -1 + \frac{1}{2}t.
\end{cases}
\]
Therefore, the complete solution set in $\mathbb{R}^4$ is:
\[
L = \left\{ \begin{pmatrix} x \\ y \\ z \\ w \end{pmatrix} = \begin{pmatrix} 1 - 2s - 2t \\ s \\ -1 + \frac{1}{2}t \\ t \end{pmatrix} = \begin{pmatrix} 1 \\ 0 \\ -1 \\ 0 \end{pmatrix} + s \begin{pmatrix} -2 \\ 1 \\ 0 \\ 0 \end{pmatrix} + t \begin{pmatrix} -2 \\ 0 \\ \frac{1}{2} \\ 1 \end{pmatrix} \;\middle|\; s, t \in \mathbb{R} \right\}.
\]
\end{exercisesolution}

% Solution by Gemini 3.7 Flash (High)
\begin{exercisesolution}[Solution to \cref{exc:simultaneous_row_reduction}]
\label{sol:simultaneous_row_reduction}
\exinfo{Solution by Gemini 3.7 Flash (High). Applies simultaneous Gaussian elimination on a dual augmented matrix.}
Since both systems $(S_1)$ and $(S_2)$ share the same coefficient matrix $A$, we form the $3 \times 5$ augmented matrix $(A \mid b_1 \mid b_2)$ and perform row reduction:
\[
\begin{pmatrix}
1 & 1 & 1 & \big| & 6 & 7 \\
1 & 2 & 2 & \big| & 11 & 10 \\
2 & 3 & -4 & \big| & 3 & 3
\end{pmatrix}
\xrightarrow{\begin{subarray}{l} -1 \cdot E_1 + E_2 \to E_2 \\ -2 \cdot E_1 + E_3 \to E_3 \end{subarray}}
\begin{pmatrix}
1 & 1 & 1 & \big| & 6 & 7 \\
0 & 1 & 1 & \big| & 5 & 3 \\
0 & 1 & -6 & \big| & -9 & -11
\end{pmatrix}
\]
Eliminate the second column in row 3:
\[
\xrightarrow{-1 \cdot E_2 + E_3 \to E_3}
\begin{pmatrix}
1 & 1 & 1 & \big| & 6 & 7 \\
0 & 1 & 1 & \big| & 5 & 3 \\
0 & 0 & -7 & \big| & -14 & -14
\end{pmatrix}
\xrightarrow{-\frac{1}{7} \cdot E_3 \to E_3}
\begin{pmatrix}
1 & 1 & 1 & \big| & 6 & 7 \\
0 & 1 & 1 & \big| & 5 & 3 \\
0 & 0 & 1 & \big| & 2 & 2
\end{pmatrix}
\]
Now clear the entries above the third pivot:
\[
\xrightarrow{\begin{subarray}{l} -1 \cdot E_3 + E_2 \to E_2 \\ -1 \cdot E_3 + E_1 \to E_1 \end{subarray}}
\begin{pmatrix}
1 & 1 & 0 & \big| & 4 & 5 \\
0 & 1 & 0 & \big| & 3 & 1 \\
0 & 0 & 1 & \big| & 2 & 2
\end{pmatrix}
\xrightarrow{-1 \cdot E_2 + E_1 \to E_1}
\begin{pmatrix}
1 & 0 & 0 & \big| & 1 & 4 \\
0 & 1 & 0 & \big| & 3 & 1 \\
0 & 0 & 1 & \big| & 2 & 2
\end{pmatrix}
\]
Reading off the solutions from the respective right-hand columns:
\begin{itemize}
    \item For $(S_1)$, the unique solution is $(x, y, z) = (1, 3, 2)$.
    \item For $(S_2)$, the unique solution is $(x, y, z) = (4, 1, 2)$.
\end{itemize}
\end{exercisesolution}

% Solution by Gemini 3.7 Flash (High)
\begin{exercisesolution}[Solution to \cref{exc:two_parameter_linear_system}]
\label{sol:two_parameter_linear_system}
\exinfo{Solution by Gemini 3.7 Flash (High). Analyzes consistency parameters $a, b$ and determines the unique solution.}
We write the augmented matrix of the system and apply Gaussian elimination:
\[
\begin{pmatrix}
1 & 2 & -3 & \big| & 10 \\
3 & 7 & -9 & \big| & b \\
-3 & -5 & a & \big| & b \\
2 & 4 & -3 & \big| & 5
\end{pmatrix}
\xrightarrow{\begin{subarray}{l} -3 \cdot E_1 + E_2 \to E_2 \\ 3 \cdot E_1 + E_3 \to E_3 \\ -2 \cdot E_1 + E_4 \to E_4 \end{subarray}}
\begin{pmatrix}
1 & 2 & -3 & \big| & 10 \\
0 & 1 & 0 & \big| & b - 30 \\
0 & 1 & a - 9 & \big| & b + 30 \\
0 & 0 & 3 & \big| & -15
\end{pmatrix}
\]
From row 4, we have $3z = -15 \implies z = -5$.
Next, perform $-1 \cdot E_2 + E_3 \to E_3$:
\[
\xrightarrow{-1 \cdot E_2 + E_3 \to E_3}
\begin{pmatrix}
1 & 2 & -3 & \big| & 10 \\
0 & 1 & 0 & \big| & b - 30 \\
0 & 0 & a - 9 & \big| & 60 \\
0 & 0 & 3 & \big| & -15
\end{pmatrix}
\]
The third row represents the equation $(a - 9)z = 60$. Substituting $z = -5$ gives:
\[
(a - 9)(-5) = 60 \implies a - 9 = -12 \implies a = -3.
\]
Therefore:
\begin{itemize}
    \item If $a \ne -3$, the third and fourth rows are contradictory, and the system has \textbf{no solutions}.
    \item If $a = -3$, the system is consistent for \textbf{every} $b \in \mathbb{R}$.
\end{itemize}
When $a = -3$, we determine the unique solution:
\begin{itemize}
    \item $z = -5$,
    \item $y = b - 30$,
    \item From row 1: $x + 2(b - 30) - 3(-5) = 10 \implies x + 2b - 60 + 15 = 10 \implies x = 55 - 2b$.
\end{itemize}
In conclusion, the system admits a solution if and only if $a = -3$ (with $b \in \mathbb{R}$ arbitrary), and the unique solution is $(x, y, z) = (55 - 2b, b - 30, -5)$.
\end{exercisesolution}

% Solution by Gemini 3.7 Flash (High)
\begin{exercisesolution}[Solution to \cref{exc:parametric_linear_system_solvability}]
\label{sol:parametric_linear_system_solvability}
\exinfo{Solution by Gemini 3.7 Flash (High). Computes the solvability hyperplane condition $b_1 - 4b_2 - b_3 = 0$ via row reduction.}
We form the augmented matrix and apply row operations:
\[
\begin{pmatrix}
0 & 7 & -2 & \big| & b_1 \\
-1 & 2 & -\frac{1}{2} & \big| & b_2 \\
4 & -1 & 0 & \big| & b_3
\end{pmatrix}
\xrightarrow{E_1 \leftrightarrow E_2}
\begin{pmatrix}
-1 & 2 & -\frac{1}{2} & \big| & b_2 \\
0 & 7 & -2 & \big| & b_1 \\
4 & -1 & 0 & \big| & b_3
\end{pmatrix}
\xrightarrow{-1 \cdot E_1 \to E_1}
\begin{pmatrix}
1 & -2 & \frac{1}{2} & \big| & -b_2 \\
0 & 7 & -2 & \big| & b_1 \\
4 & -1 & 0 & \big| & b_3
\end{pmatrix}
\]
Eliminate the first column of row 3:
\[
\xrightarrow{-4 \cdot E_1 + E_3 \to E_3}
\begin{pmatrix}
1 & -2 & \frac{1}{2} & \big| & -b_2 \\
0 & 7 & -2 & \big| & b_1 \\
0 & 7 & -2 & \big| & b_3 + 4b_2
\end{pmatrix}
\xrightarrow{-1 \cdot E_2 + E_3 \to E_3}
\begin{pmatrix}
1 & -2 & \frac{1}{2} & \big| & -b_2 \\
0 & 7 & -2 & \big| & b_1 \\
0 & 0 & 0 & \big| & b_3 + 4b_2 - b_1
\end{pmatrix}
\]
\textbf{Consistency condition:} The system possesses at least one solution if and only if the third row does not produce a contradiction ($0 = c$ with $c \ne 0$). Hence, the necessary and sufficient condition is:
\[
b_1 - 4b_2 - b_3 = 0.
\]
\textbf{Solution set when consistent:} When $b_1 - 4b_2 - b_3 = 0$, $z$ is a free variable. Setting $z = 14t$ for $t \in \mathbb{R}$:
\[
7y - 2z = b_1 \implies 7y = b_1 + 28t \implies y = \frac{1}{7}b_1 + 4t.
\]
From the first row:
\[
x = -b_2 + 2y - \frac{1}{2}z = -b_2 + 2\left(\frac{1}{7}b_1 + 4t\right) - 7t = \frac{2}{7}b_1 - b_2 + t.
\]
Thus, the solution set $S \subseteq \mathbb{R}^3$ is:
\[
S = \left\{ \begin{pmatrix} \frac{2}{7}b_1 - b_2 \\ \frac{1}{7}b_1 \\ 0 \end{pmatrix} + t \begin{pmatrix} 1 \\ 4 \\ 14 \end{pmatrix} \;\middle|\; t \in \mathbb{R} \right\}.
\]
\end{exercisesolution}

% Solution by Gemini 3.7 Flash (High)
\begin{exercisesolution}[Solution to \cref{exc:inconsistent_2x2_system}]
\label{sol:inconsistent_2x2_system}
\exinfo{Hybrid Solution (Gemini 3.7 Flash / Official Problem Sheet 4, Exercise 1). Proves the contrapositive geometric condition for $2 \times 2$ inconsistent systems using both determinant invertibility and echelon reduction.}
Let $A = \begin{pmatrix} a_{11} & a_{12} \\ a_{21} & a_{22} \end{pmatrix} \in \M_{2 \times 2}(\mathbb{R})$ and $b = \begin{pmatrix} b_1 \\ b_2 \end{pmatrix} \in \mathbb{R}^2$. 

Suppose for contradiction that neither row of $A$ is a scalar multiple of the other. This means that the row vectors $(a_{11}, a_{12})$ and $(a_{21}, a_{22})$ are linearly independent in $\mathbb{R}^2$. 

For a $2 \times 2$ matrix, linear independence of the rows is equivalent to
\[
\det(A) = a_{11}a_{22} - a_{12}a_{21} \ne 0.
\]
When $\det(A) \ne 0$, the matrix $A$ is invertible with inverse
\[
A^{-1} = \frac{1}{\det(A)} \begin{pmatrix} a_{22} & -a_{12} \\ -a_{21} & a_{11} \end{pmatrix}.
\]
Therefore, the system $Ax = b$ admits the unique solution $x = A^{-1}b$. But this contradicts the hypothesis that the system $Ax = b$ has no solutions.

Hence, the rows of $A$ must be linearly dependent, which means that one row of $A$ is a scalar multiple of the other.

\begin{ainote}
\textbf{Alternative Elementary Proof via Row Operations:}
We can also prove the contrapositive directly using only Gaussian elimination: suppose no row of $A = \begin{pmatrix} a & b \\ c & d \end{pmatrix}$ is a scalar multiple of the other.
\begin{itemize}
    \item If $a \ne 0$, subtract $\frac{c}{a} R_1$ from $R_2$ to obtain the second row $\bigl(0, \; d - \frac{bc}{a} \;\big|\; b_2 - \frac{c}{a}b_1\bigr)$. Since the rows are not proportional, $ad - bc \ne 0$, so the pivot $d - \frac{bc}{a} \ne 0$, which allows solving for both unknowns uniquely for any $(b_1, b_2)$.
    \item If $a = 0$, then $c \ne 0$ (otherwise row 1 is $0$ times row 2). Swapping rows reduces to the previous case.
\end{itemize}
Thus, if no row is a multiple of the other, the system is consistent for every right-hand side $b \in \mathbb{R}^2$.
\end{ainote}
\end{exercisesolution}
